\cotefile{93}
\scanpage{3}
\margin{\handwritten{G? = correction de Grothendieck}}
\margin{\handwritten{Pour une version un brin plus \unsure{complete|complète} : Lecture Notes 476 (Anvers $\overline{\underline{\mathrm{IV}}}$) p 53--73 see also : Tate, ibidem p 33--52 et \unsure{spécialement} p 33--39}}

\begin{center}
COURBES ELLIPTIQUES : FORMULAIRE

d'après J. Tate\\
mis au goût du jour par P. Deligne
\end{center}

1.\qquad \uline{Formules générales}.

Dans ce formulaire, on appellera courbe de genre 1 sur un schéma de base $S$ un morphisme propre et plat de présentation finie
\[
p : E \handwritten{\longrightarrow} S \qquad ,
\]
muni d'une section $e$ contenue dans l'ouvert de lissité, dont les fibres géométriques sont des courbes réduites irréductibles de genre \margin{\handwritten{G?}}\handwritten{\uline{ar}}ithmétique 1 . Les fibres géométriques de $p$ sont donc soit :

a) une courbe elliptique, i.e. une courbe propre, lisse et connexe de genre 1 ,

b) une droite projective dont on a identifié deux points distincts (cubique dans $P_2$ à point double ordinaire),

c) une droite projective dont on a identifié deux points infiniment voisins (cubique dans $P_2$ à singularité cuspidale).

On dira dans le cas a) (resp. b), c)) que la fibre est de type elliptique (resp. multiplicatif, resp. additif).

Si $E$ est une courbe de genre 1 sur $S$, on définit un faisceau inversible $\omega$ sur $S$ par la formule
\[
\omega=e^*\Omega^1_{E/S}\qquad .
\]
La section $e$ est un diviseur de Cortier relatif de $E$ sur $S$ ;

% archive sheet number: 2
